\documentclass[12pt]{article}
\usepackage{geometry}
\geometry{a4paper, margin=1in}
\usepackage{amsmath}
\usepackage{amsfonts}
\usepackage{setspace}
\usepackage{parskip}
\usepackage{enumitem}
\usepackage{xeCJK}
\setCJKmainfont{Noto Serif CJK SC}
\usepackage{booktabs}
\usepackage{longtable}
\usepackage[utf8]{inputenc}
\usepackage[T1]{fontenc}
\usepackage{tocloft}
\usepackage{hyperref}
\hypersetup{colorlinks=true, linkcolor=blue, urlcolor=blue}
\usepackage{natbib}
\usepackage{authblk}

\setlength{\parindent}{0pt}
\setlength{\parskip}{1em}
\onehalfspacing

\begin{document}

\title{研究計畫書：AI輔助製作中文化NEJM病例教學講義對住院醫師學習效果之影響}
\author{新竹台大分院教學部 朱秀宸/謝慕揚}
\date{2025年7月19日}
\maketitle

\begin{abstract}
本研究旨在探討AI輔助製作中文化NEJM病例教學講義對住院醫師學習效果的影響，採用前後對照與混合方法研究設計。研究對象包括住院醫師（R1-R3）、主治醫師及實習醫師，比較AI輔助中文化教材與傳統英文或中文教材的學習成效。主要結果指標包括病例理解測驗分數、學習效率及知識保留率；次要指標包括學習滿意度、參與度及自信心。研究分三階段執行：準備期（1-2個月）、介入期（3-4個月）及評估期（1-2個月）。預計每組32人，採用配對t檢定、獨立樣本t檢定及重複測量ANOVA進行分析，質性資料以主題分析法處理。研究結果將投稿至《Medical Teacher》、《BMC Medical Education》等期刊，期為華語醫學教育提供創新且可複製的模式。
\end{abstract}

\section{研究背景與目的}
隨著人工智慧（AI）技術的進步，其在醫學教育中的應用日益受到重視。NEJM病例報告以其高品質的臨床案例聞名，但英文內容對非英語母語的住院醫師可能構成障礙。本研究旨在評估AI輔助製作的中文化NEJM病例講義是否能提升住院醫師的學習效果，並探討其本土化教材製作的可行性，為華語醫學教育提供實證基礎。

\subsection{研究問題}
\begin{enumerate}
    \item AI輔助中文化NEJM病例講義是否能提升住院醫師的病例理解、學習效率及知識保留？
    \item AI輔助教材與傳統教材相比，在學習滿意度、參與度及自信心方面有何差異？
    \item AI輔助教材製作的優缺點及改進方向為何？
\end{enumerate}

\subsection{研究目標}
\begin{itemize}
    \item 比較AI輔助中文化教材與傳統教材在學習效果上的差異。
    \item 評估AI輔助教材的臨床相關性與實用性。
    \item 探索AI技術在醫學教育本土化中的潛力與限制。
\end{itemize}

\section{研究設計}
\subsection{研究類型}
\begin{itemize}
    \item 前後對照研究：比較介入前後學習效果的變化。
    \item 混合方法研究：結合量化（問卷與測驗）與質性（訪談）資料。
\end{itemize}

\subsection{研究對象}
\begin{itemize}
    \item 住院醫師（R1-R3）
    \item 主治醫師（作為輔助評估者）
    \item 醫學生（實習醫師，作為次要對象）
\end{itemize}

\subsection{介入措施}
\begin{itemize}
    \item 實驗組：使用AI輔助製作的中文化NEJM病例講義。
    \item 對照組：使用傳統英文NEJM病例或傳統中文教材。
\end{itemize}

\subsection{測量指標}
\subsubsection{主要結果指標}
\begin{itemize}
    \item 病例理解測驗分數（客觀結構化試題）。
    \item 學習效率（完成學習所需時間）。
    \item 知識保留率（1個月後追蹤測驗）。
\end{itemize}

\subsubsection{次要結果指標}
\begin{itemize}
    \item 學習滿意度（5點Likert量表）。
    \item 參與度（討論參與頻率、提問次數）。
    \item 自信心（病例診斷自信度評分）。
\end{itemize}

\section{研究方法}
\subsection{問卷設計}
\subsubsection{量化部分（5點Likert量表）}
\begin{itemize}
    \item 內容品質
    \begin{itemize}
        \item 中文翻譯準確性
        \item 醫學術語適切性
        \item 內容組織邏輯性
    \end{itemize}
    \item 學習體驗
    \begin{itemize}
        \item 理解容易度
        \item 學習動機提升
        \item 記憶深刻程度
    \end{itemize}
    \item 實用性
    \begin{itemize}
        \item 臨床相關性
        \item 應用便利性
        \item 時間節省效果
    \end{itemize}
\end{itemize}

\subsubsection{質性部分（開放式問題）}
\begin{itemize}
    \item 與原版英文NEJM相比的優缺點
    \item 對AI輔助教材製作的看法
    \item 改進建議
    \item 未來使用意願
\end{itemize}

\subsection{研究執行流程}
\begin{longtable}{p{3cm} p{4cm} p{8cm}}
    \toprule
    \textbf{階段} & \textbf{時間} & \textbf{工作內容} \\
    \midrule
    第一階段：準備期 & 第1-2個月 & 
    \begin{itemize}[noitemsep]
        \item 完成10份AI輔助中文化NEJM病例講義。
        \item 設計前測問卷及測驗題目。
        \item 提交並獲得IRB倫理審查通過。
        \item 招募研究對象，隨機分組。
    \end{itemize} \\
    第二階段：介入期 & 第3-6個月 & 
    \begin{itemize}[noitemsep]
        \item 進行前測評估（基線測驗及問卷）。
        \item 實驗組使用AI輔助教材，對照組使用傳統教材。
        \item 實施分組教學（每週1次，共12週）。
        \item 收集即時回饋（課後問卷及討論記錄）。
    \end{itemize} \\
    第三階段：評估期 & 第7-8個月 & 
    \begin{itemize}[noitemsep]
        \item 進行後測評估（測驗及問卷）。
        \item 對10-15名參與者進行深度訪談。
        \item 實施1個月後追蹤測驗。
        \item 整理數據，進行統計分析。
    \end{itemize} \\
    \bottomrule
\end{longtable}

\subsection{時程表}
\begin{longtable}{p{2cm} p{4cm} p{4cm} p{4cm}}
    \toprule
    \textbf{時間} & \textbf{階段} & \textbf{任務} & \textbf{負責人} \\
    \midrule
    第1個月 & 準備期 & 製作AI輔助教材 & 研究助理 \\
    第1-2個月 & 準備期 & 設計問卷與測驗 & 研究團隊 \\
    第2個月 & 準備期 & IRB申請與審查 & 主研究者 \\
    第2個月 & 準備期 & 招募與分組 & 研究助理 \\
    第3-6個月 & 介入期 & 前測與分組教學 & 教學團隊 \\
    第3-6個月 & 介入期 & 即時回饋收集 & 研究助理 \\
    第7個月 & 評估期 & 後測與訪談 & 研究團隊 \\
    第8個月 & 評估期 & 追蹤測驗 & 研究助理 \\
    第8個月 & 評估期 & 數據分析與報告撰寫 & 統計分析師 \\
    \bottomrule
\end{longtable}

\section{統計分析計畫}
\subsection{樣本數計算}
\begin{itemize}
    \item 預期效果量：中等效果（Cohen's d = 0.5）。
    \item 檢定力：80\%。
    \item 顯著水準：0.05。
    \item 建議樣本數：每組至少32人（共64人）。
\end{itemize}

\subsection{分析方法}
\begin{itemize}
    \item 配對t檢定：比較前後測分數變化。
    \item 獨立樣本t檢定：比較實驗組與對照組。
    \item 重複測量ANOVA：分析追蹤測驗效果。
    \item 質性分析：使用主題分析法處理訪談資料。
\end{itemize}

\section{預期成果}
\begin{itemize}
    \item 驗證AI輔助中文化教材在學習效果上的優勢。
    \item 提供可複製的醫學教育本土化模式。
    \item 為AI技術在醫學教育的應用提供實證支持。
\end{itemize}

\section{期刊投稿計畫}
\subsection{目標期刊}
\begin{itemize}
    \item \textit{Medical Teacher} (IF: 3.3)
    \item \textit{BMC Medical Education} (IF: 2.9)
    \item \textit{Academic Medicine} (IF: 7.9)
    \item 中華醫學教育雜誌
    \item 台灣醫學教育學刊
\end{itemize}

\subsection{投稿角度}
\begin{itemize}
    \item 創新性：首次系統性評估AI輔助醫學教育材料。
    \item 實用性：提出可複製的教育改善模式。
    \item 本土化：滿足華語醫學教育需求。
\end{itemize}

\section{研究限制與對策}
\subsection{潛在限制}
\begin{itemize}
    \item 樣本偏差：單一醫院可能限制結果外推性。
    \item 新奇效應：AI教材可能因新穎性影響滿意度。
    \item 語言能力差異：參與者的英文能力可能影響對照組表現。
\end{itemize}

\subsection{對策}
\begin{itemize}
    \item 多中心研究：與多家醫院合作，增加樣本多樣性。
    \item 延長觀察期：降低新奇效應影響。
    \item 控制基線特徵：記錄並分析參與者的語言能力。
\end{itemize}

\section{後續發展建議}
\begin{itemize}
    \item 擴大規模：納入多醫院、多科別參與者。
    \item 技術改進：優化AI翻譯與教材生成演算法。
    \item 效果追蹤：評估長期學習成效。
    \item 成本效益：分析AI教材製作的資源投入回報。
\end{itemize}

\section{參考文獻}
\bibliographystyle{plain}
\bibliography{references}


\section{附件- 問卷}
感謝您參與本研究！本問卷旨在評估AI輔助製作的中文化NEJM病例教學講義對學習效果的影響。問卷包含兩部分：\textbf{量化部分}（使用5點Likert量表）與\textbf{質性部分}（開放式問題）。請依據您的實際體驗如實回答，所有資料將匿名處理，僅用於學術研究。完成問卷約需10-15分鐘。

\subsection*{量化部分：5點Likert量表}
請根據您的體驗，對以下陳述進行評分：\\
1 = 完全不同意，2 = 不同意，3 = 無意見，4 = 同意，5 = 完全同意

\begin{longtable}{p{1cm} p{7cm} p{1.5cm} p{1.5cm} p{1.5cm} p{1.5cm} p{1.5cm}}
\toprule
\textbf{編號} & \textbf{陳述} & \textbf{1} & \textbf{2} & \textbf{3} & \textbf{4} & \textbf{5} \\
\midrule
\multicolumn{7}{l}{\textbf{內容品質}} \\
1 & 本教材的中文翻譯準確且易於理解。 & & & & & \\
2 & 本教材使用的醫學術語適切且符合臨床實務。 & & & & & \\
3 & 本教材的內容組織邏輯清晰，結構合理。 & & & & & \\
\multicolumn{7}{l}{\textbf{學習體驗}} \\
4 & 本教材有助於我更容易理解病例內容。 & & & & & \\
5 & 本教材提升了我的學習動機。 & & & & & \\
6 & 本教材的內容讓我記憶更深刻。 & & & & & \\
\multicolumn{7}{l}{\textbf{實用性}} \\
7 & 本教材的內容與臨床實務高度相關。 & & & & & \\
8 & 本教材的設計便於在學習中應用。 & & & & & \\
9 & 本教材有助於節省我的學習時間。 & & & & & \\
\multicolumn{7}{l}{\textbf{學習滿意度與自信心}} \\
10 & 我對使用本教材的整體學習體驗感到滿意。 & & & & & \\
11 & 使用本教材後，我對病例診斷的自信心有所提升。 & & & & & \\
12 & 我認為本教材有助於提升我的參與度和討論積極性。 & & & & & \\
\bottomrule
\caption{量化問卷：請在對應評分欄中打勾}
\end{longtable}

\subsection*{質性部分：開放式問題}
請詳細回答以下問題，您的意見對我們改進教材至關重要。

\begin{longtable}{p{1cm} p{14cm}}
\toprule
\textbf{編號} & \textbf{問題} \\
\midrule
1 & 與原版英文NEJM病例報告相比，您認為AI輔助中文化教材的優點與缺點分別是什麼？ \\
 & \vspace{3cm} \\
2 & 您對AI輔助製作醫學教育教材的看法如何？ \\
 & \vspace{3cm} \\
3 & 您對本教材有何改進建議？ \\
 & \vspace{3cm} \\
4 & 您是否願意在未來繼續使用AI輔助中文化教材？請說明原因。 \\
 & \vspace{3cm} \\
\bottomrule
\caption{質性問卷：請在空格中詳細書寫您的回答}
\end{longtable}

\subsection*{參與度觀察表（由教學者填寫）}
請教學者根據觀察，記錄每位參與者在教學過程中的表現：

\begin{longtable}{p{8cm} p{8cm}}
\toprule
\textbf{觀察項目} & \textbf{記錄} \\
\midrule
討論參與頻率（次數/每場教學） & \\
提問次數（次數/每場教學） & \\
其他觀察（例如主動性、互動程度等） & \vspace{2cm} \\
\bottomrule
\caption{參與度觀察表}
\end{longtable}

\subsection*{病例理解測驗（附錄）}
\begin{itemize}
    \item 本部分將於教學結束後進行，包含10道客觀結構化試題，涵蓋病例診斷、治療方案及相關知識。
    \item 測驗時間：30分鐘。
    \item 評分標準：每題10分，總分100分。
\end{itemize}

\subsection*{追蹤測驗（1個月後）}
\begin{itemize}
    \item 內容：與前測相同，重新評估知識保留率。
    \item 時間：30分鐘。
    \item 目的：測量學習效果的長期維持。
\end{itemize}

\subsection*{注意事項}
\begin{itemize}
    \item 請確保所有問題均已回答，特別是開放式問題請提供具體意見。
    \item 問卷將於每次教學結束後及追蹤階段收集。
    \item 如有任何問題，請聯繫研究團隊。
\end{itemize}

\end{document}